%!TEX program = lualatex
\documentclass[11pt, a4paper]{scrartcl}
% Packages
\usepackage[margin=1.25in]{geometry}
\usepackage{index}
\usepackage{amsbsy} % Bold math symbols
\makeindex
\usepackage{tcolorbox}
\usepackage{varwidth}
\usepackage{amsmath, amssymb}
\usepackage{esint}
\usepackage{titlesec}
\usepackage{xcolor}
\usepackage{titling}
\usepackage[linktocpage]{hyperref}
\usepackage{pgfplots}
\usepackage{multicol}
\setlength{\columnsep}{2em}
\usepackage{caption}
\usepackage{amsthm}
\usepackage{import}
\usepackage{cancel}
\usepackage{caption}
\usepackage{nicematrix}
\usepackage{mathtools}
\usepackage{enumerate}
\usepackage{graphicx}
\usepackage{makecell}
\usepackage{tikz}
\usepackage{circuitikz}
\usepackage{pgf-spectra}
\usepackage[version=4]{mhchem}
\usetikzlibrary{arrows.meta,calc,positioning, angles,quotes,patterns.meta,intersections}
\usepackage[italian]{babel}
% Font
%\usepackage{fontspec}
%\setmainfont{Libertinus Serif}

%\usepackage{unicode-math}
%\setmathfont{Libertinus Math}


% To reset footnote numbering each page
\usepackage[perpage]{footmisc}
\usepackage{setspace}
\setstretch{1.2}

% Titles 
\title{Note di Cosmologia}
\author{Manuel Deodato}
\date{}


% svolgimento
\newenvironment{svolgimento}{\renewcommand\qedsymbol{$\blacksquare$}\begin{proof}[Svolgimento]}{\end{proof}}

\newtheoremstyle{style}% name of the style to be used
{5pt}% measure of space to leave above the theorem. E.g.: 3pt
{5pt}% measure of space to leave below the theorem. E.g.: 3pt
{\normalfont}% name of font to use in the body of the theorem
%{15pt}% measure of space to indent
{0pt}% measure of space to indent
{\noindent\bfseries}% name of head font
{}% punctuation between head and body
{ }% space after theorem head; " " = normal interword space
{\thmname{#1}\thmnumber{ #2}{\thmnote{ (#3)}.\ }}


\theoremstyle{style}
\newtheorem{esempio}{Esempio}[section]
\newtheorem{definizione}{Definizione}[section]
\newtheorem{prop}{Proposizione}[section]
\newtheorem{teorema}{Teorema}[section]
\newtheorem{lemma}{Lemma}[teorema]
\newtheorem{corollario}{Corollario}[teorema]
\newtheorem{osservazione}{Osservazione}[section]
\newtheorem{assunzione}{Assunzione}[section]
\newtheorem{notazione}{Notazione}[section]

\newtheoremstyle{style1}% name of the style to be used
{5pt}% measure of space to leave above the theorem. E.g.: 3pt
{5pt}% measure of space to leave below the theorem. E.g.: 3pt
{\normalfont}% name of font to use in the body of the theorem
%{15pt}% measure of space to indent
{0pt}% measure of space to indent
{\noindent\bfseries\color{red}}% name of head font
{}% punctuation between head and body
{ }% space after theorem head; " " = normal interword space
{\thmname{#1}\thmnumber{ #2}{\thmnote{ (#3)}.\ }}

\theoremstyle{style1}
\newtheorem{esercizio}{Esercizio}[section]


%% Generic box
\newtcolorbox{eqbox}[1][]
{
colback=gray!10,
arc=0pt,
boxrule=0pt,
title=#1
}

 \newenvironment{boxenv}[1][]{
    \begin{eqbox}[#1]
    }{
   \end{eqbox}
}



%%%%%%%%%% Medie con integrali multipli
\def\Yint#1{\mathchoice
    {\YYint\displaystyle\textstyle{#1}}%
    {\YYint\textstyle\scriptstyle{#1}}%
    {\YYint\scriptstyle\scriptscriptstyle{#1}}%
    {\YYint\scriptscriptstyle\scriptscriptstyle{#1}}%
      \!\iint}
\def\YYint#1#2#3{{\setbox0=\hbox{$#1{#2#3}{\iint}$}
    \vcenter{\hbox{$#2#3$}}\kern-.51\wd0}}
\def\longdash{{-}\mkern-3.5mu{-}} 
   % consider using "\mkern-7.5mu" if esint package is loaded
\def\tiltlongdash{\rotatebox[origin=c]{15}{$\longdash$}}
\def\fiint{\Yint\tiltlongdash}

\def\Zint#1{\mathchoice
    {\YYint\displaystyle\textstyle{#1}}%
    {\YYint\textstyle\scriptstyle{#1}}%
    {\YYint\scriptstyle\scriptscriptstyle{#1}}%
    {\YYint\scriptscriptstyle\scriptscriptstyle{#1}}%
      \!\iiint}
      \def\tilongdash{\mkern6mu{-}\mkern-4mu{-}\mkern-5mu{-}} 
   % consider using "\mkern-7.5mu" if esint package is loaded
\def\titiltlongdash{\rotatebox[origin=c]{15}{$\tilongdash$}}
\def\fiiint{\Zint\titiltlongdash}

%Captions
\captionsetup[figure]{font=footnotesize,labelfont=footnotesize}
\captionsetup[table]{font=footnotesize,labelfont=footnotesize}
%Titlesec
\titleformat{\section}
{\fontsize{15}{20}\sffamily\scshape}
{\normalfont\color{gray}{\fontsize{20}{20}\selectfont\thesection}}
{0.7em}
{}
\hypersetup{colorlinks,breaklinks, linkcolor=[RGB]{74, 122, 164}}
\definecolor{asdf}{HTML}{4a7aa4}
% Personalizza la formattazione della subsection
\titleformat{\subsection}[block]{\fontsize{13}{20}\bfseries}{\normalfont\thesubsection}{.5em}{}


% Personalizza la formattazione della subsubsection
\titleformat{\subsubsection}[block]{\fontsize{12}{20}\bfseries}{\normalfont\thesubsubsection}{.5em}{}

% Maketitle customization
\renewcommand{\maketitle}{
\begin{center}
{\sffamily
{\fontsize{20}{20}\selectfont\MakeUppercase\thetitle}}

\vspace{0.2in}

{\large\scshape\sffamily\theauthor}
\end{center}
}


% Differenziale 
\newcommand{\dd}{\mathop{}\!\mathrm{d}}




%Evaluate symbol
\DeclareMathOperator{\di}{d\!}
\newcommand*\Eval[3]{\left.#1\right\rvert_{#2}^{#3}}

%%%%%%% Numero delle equazioni in formato a.b
\numberwithin{equation}{subsection}
%%%%%

%%%%%%%%%% Personalizzazione numeri lista
\renewcommand{\theenumi}{(\arabic{enumi})}

%%%% Table of contents

\usepackage[titles]{tocloft}

\renewcommand{\cftdot}{}
\usepackage{titletoc}
%\setcounter{tocdepth}{2}

%%%%%%%%%%%%%%%% Toc style

% Personalizzazione scritta indice


%%%% Lambda bar
\newcommand{\lbar}{{\mkern0.75mu\mathchar '26\mkern -9.75mu\lambda}}


\begin{document}
\maketitle
\newpage
\tableofcontents 
\newpage
\section{Introduzione}
\subsection{Sviluppi storici in cosmologia}
\subsubsection{Paradosso di Olbers (1826)}
Il contesto storico in cui si presenta questa trattazione vede l'universo come infinito, eterno, statico, omogeneo, isotropo; ci si riferir\`a a questo modello come \textit{cosmologia newtoniana}, che rappresenta lo studio dell'universo prima della relativit\`a generale, quindi facendo uso della gravit\`a proposta da Newton.

Il paradosso di Olbers, noto anche come \textit{paradosso del cielo stellato}, evidenzia come se l'universo presentasse le caratteristiche evidenziate sopra, allora lungo ogni direzione nel cielo si deve poter osservare almeno una stella; per questa ragione, il cielo dovrebbe apparire completamente luminoso.
Il flusso di luce ricevuta sulla superficie terrestre \`e dato da
\[
\Phi_s = \int _0^{+\infty} \frac{L_\odot }{4 \pi r^2} n_s 4 \pi r^2 \dd r
\] 
con $n_s$ densit\`a delle stelle e $L_\odot $ luminosit\`a stellare (posta per semplicit\`a pari a quella del Sole).
In questo integrale, $\dd N = n_s 4 \pi r^2 \dd r$ rappresenta il numero di stelle contenute in un guscio sferico di raggio $r$ e spessore $\dd r$.
Per assunzione, $n_s$ deve essere uniforme, quindi
\[
\Phi_s = L_\odot  n_s \int _0^{+\infty} \dd r = +\infty
\] 
L'\textit{ipotesi di Olbers} si propone di risolvere questo problema assumendo l'esistenza di polvere nel mezzo interstellare/intergalattico, che causa l'assorbimento della radiazione proveniente dalle stelle.
Questo, per\`o, non pu\`o funzionare: una nube che assorbe la radiazione tende, per tempi sufficientemente lunghi, a scaldarsi ed emettere come un corpo nero alla stessa temperatura delle stelle, finendo per avere lo stesso identico problema di partenza.

\subsubsection{Instabilit\`a gravitazionale dell'universo di Newton}

Il problema della stabilit\`a si basa su come risponde l'universo newtoniano ad una piccola perturbazione di densit\`a.
Partendo dal vedere l'universo come contenente materia distribuita uniformemente al suo interno, l'idea della sua stabilit\`a si fonda sul fatto che, in presenza di una distribuzione completamente omogenea, non esiste una direzione privilegiata e una configurazione statica pu\`o essere considerata in equilibrio ideale.
Tuttavia, questa situazione di equilibrio \`e solamente ideale: immaginando che, per qualunque ragione, una regione diventi leggermente pi\`u densa della media, questa eserciterebbe un'attrazione maggiore sulle zone circostanti, richiamando altra materia.
Quest'accrescimento nella regione in questione e la rispettiva diminuzione delle regioni da cui la materia viene sottratta.
Si evidenzia, allora, come il sistema non smorza piccole perturbazioni di densit\`a, ma le alimenta, risultando, di conseguenza, instabile.

Dal punto di vista dinamico, l'equazione del moto per un guscio sferico di raggio $R(t)$ contenente una massa $M$ \`e
\[
\ddot{R} = - \frac{GM}{R^2}
\] 
Se il sistema \`e inizialmente fermo, allora $\dot{R} = 0$, per\`o si vede che $\ddot{R} < 0$, per cui il raggio inizia a diminuire e si verifica un collasso, il che alimenta le perturbazioni iniziali nella densit\`a.
Su tempi scala sufficientemente lunghi, quindi, l'universo newtoniano non pu\`o esistere, in quanto instabile.

Nel 1895, Seeliger ipotizza che \textit{la gravit\`a su grandi distanze venga schermata}, cio\`e propone che il potenziale gravitazionale diventi
\[
\Phi(r)  = G \int \frac{\rho}{r} \dd ^3 r \longrightarrow G \int \frac{\rho }{r} e^{-\lambda r}\dd ^3r
\] 
Sotto questa ipotesi, l'equazione di Poisson $\nabla ^2 \Phi = 4 \pi G \rho  $ diventa 
\[
 \nabla ^2 \Phi - \lambda \Phi = 4 \pi G \rho 
\] 
Questa idea \`e poi ripresa da Einstein nel suo modello cosmologico, prendendo il nome di \textit{costante cosmologia}.

\subsubsection{Cosmologia relativistica}

La cosmologia relativistica nasce dalla relativit\`a generale e viene introdotta, per la prima volta, in un articolo di Einstein del 1917.
Questo articolo sulla cosmologia nasce dalla necessit\`a di risolvere problemi intrinseci alla relativit\`a generale. 
I principi e risultati generali di questa teoria sono i seguenti.
\begin{itemize}
	\item \textit{Principio di relativit\`a}: le leggi della fisica hanno la stessa forma in tutti i sistemi di riferimento, anche non-inerziali, purch\'e siano scritte in forma covariante.
	\item \textit{Principio di equivalenza}: localmente, gli effetti di un campo gravitazionale uniforme sono indistinguibili da quelli di un sistema accelerato; da questo, segue l'equivalenza tra massa inerziale e massa gravitazionale.
	\item \textit{Principio di covarianza generale}: le leggi della fisica devono mantenere la stessa forma sotto qualsiasi trasformazione di coordinate.
		Questo comporta che non esistono sistemi di riferimento privilegiati e che la descrizione fisica non dipende dalle coordinate scelte.
	\item \textit{Gravit\`a come geometria.} 
		La gravit\`a non \`e una forza: la presenza di massa ed energia curva lo spaziotempo e le particelle libere seguono le geodetiche di tale geometria.
		In sostanza, la materia curva lo spaziotempo e la geometria dello spaziotempo caratterizza il moto della materia.
	\item \textit{Principio di corrispondenza:} la relativit\`a generale si riduce alla gravitazione di Newton per campo gravitazionale debole e velocit\`a inferiori a quelle della luce.
\end{itemize}
L'universo di Einstein ha, come caratteristiche principali, \textit{staticit\`a} e \textit{chiusura}.
Tuttavia, l'universo descritto dall'equazione di Einstein
\[
G_{\mu \nu }  = 8 \pi G T_{\mu \nu } 
\] 
con $G_{\mu \nu } $ tensore di Einstein e $T_{\mu \nu } $ tensore energia-impulso, non descrive un universo statico, ma uno dinamico.
La dinamicit\`a, ai tempi, non era una propriet\`a nota e non era desiderata, per cui Einstein, ispirandosi al metodo proposto da Seelinger, introduce un termine di schermaggio:
\[
G_{\mu \nu }  - \Lambda g_{\mu \nu }  = 8 \pi G T_{\mu  \nu } 
\] 
dove $g_{\mu \nu } $ \`e la metrica dello spazio tempo, mentre $\Lambda $ \`e la \textit{costante cosmologica}.

Questo nuovo modello risolve il paradosso di Olbers, in quanto le equazioni di Einstein, sotto assunzione di universo statico e pieno di materia, producono un universo a curvatura positiva, quindi chiuso e, dunque, di volume finito, eliminando la divergenza nell'integrale: l'estremo superiore non \`e pi\`u $+\infty$.
Tuttavia, si ha ancora il problema dell'instabilit\`a.
L'equazione originale descriveva un universo in cui non si aveva staticit\`a: o si aveva espansione, o si aveva un collasso.
La costante cosmologica produce un effetto repulsivo su grande scala: la materia tende a formare collassi; la costante cosmologica produce espansioni.
Un opportuno valore di $\Lambda $ permette di ottenere staticit\`a.
Tuttavia, il problema della stabilit\`a rimane: uno squilibrio di densit\`a, che sia un aumento, o una diminuzioneo, comporta, rispettivamente, un collasso, oppure un'espansione (quest'ultima, legata a $\Lambda $).

Nel 1918, Shr\"odinger ipotizza che il termine cosmologico possa essere interpretato come il tensore energia-impulso di un fluido a pressione $P < 0$; in questo modo, la costante cosmologica non \`e pi\`u una forzatura matematica, ma acquisisce un significato ben preciso e si possono riscrivere le equazioni di Einstein come
\[
	G_{\mu  \nu  }  = 8 \pi G T_{\mu  \nu  }  \hspace{1cm} T_{\mu \nu }  = T_{\mu \nu } ^\text{materia} + T_{\mu \nu }^\text{(P>0)}
\] 
\subsubsection{Evidenze dell'espansione cosmologica}
Le prime osservazioni delle alte velocit\`a delle galassie esterne furono ottenute da Slipher nel 1917 tramite analisi spettrale delle righe di emissione.

La successiva determinazione delle distanze da tali galassie \`e stata ottenuta tramite il metodo delle cefeidi variabili, visto che, per distanze cos\`i grandi, il metodo della parallasse non funziona bene.
Questo metodo si basa sul fatto che la magnitudine assoluta di queste stelle \`e dipendente esclusivamente dal relativo periodo di pulsazione tramite la relazione
\[
M = a \log P + b
\] 
dove $a$ e $b$ sono parametri da calibrare sperimentalmente.
Leavitt, infatti, osserv\`o moltissime cefeidi nella piccola nube di Magellano e not\`o che quelle con periodo maggiore erano sempre le pi\`u luminose; inoltre, trovandosi a circa la stessa distanza dalla Terra, dovevano esserci differenze solamente in termini di luminosit\`a intrinseca. 
In questo modo, si pot\'e ricavare il termine $a$ confrontando magnitudine assoluta e periodo; il termine noto $b$, invece, si determina utilizzando le distanze delle cefeidi pi\`u vicine tramite il metodo della parallasse: dalla distanza e dalla magnitudine apparente $m$, si ottiene 
\[
M = m - 5 \log \left(\frac{d}{10 \text{ pc}}\right) 
\] 
Potendo determinare la magnitudine assoluta delle cefeidi nelle galassie di cui si vuole ricavare la distanza dalla Terra, quest'ultima si ottiene tramite
\[
\Phi _s = \frac{L}{4 \pi d^2}
\] 
visto che dalla magnitudine assoluta si ricava $L$.

Nel 1929, Hubble ricava, dall'osservazione di numerose galassie tramite il metodo delle cefeidi variabili, la relazione
\begin{equation}
	v = H_0d
\end{equation}
con $H_0 \approx 70 \text{ km s}^{-1} \text{ Mpc}$ costante di Hubble.
In realt\`a, la velocit\`a $v$ non \`e propriamente un osservabile, in quanto si misura il red-shift; sar\`a opportuno, quindi, esprimere la legge di Hubble in termini del red-shift.

\subsubsection{Universo di Einstein-de Sitter}

A seguito della scoperta dell'espansione dell'universo tramite la legge di Hubble, Einstein rivaluta la forma originale delle sue equazioni senza costante cosmologica e si introduce il modello di universo di Einstein-de Sitter.
Le sue caratteristiche sono una curvatura $k=0$ (universo piatto), $\Lambda =0$ e dominato esclusivamente da materia (polvere) non-relativistica a pressione trascurabile $P\simeq 0$.
\subsection{Fondamenti di relativit\`a generale}
















\end{document}
